\documentclass[../../main.tex]{subfiles}

\graphicspath{{../../image/}{../../figure/}}

\begin{document}

\section{习题}\label{section:ch09-习题}

\begin{exercise}\label{exercise:ch09-01}
（★）通过依次考虑\cref{subsection:ch09-对称性不变性}介绍的 $4$ 个公理，证明以复合为运算，绕正方形中心旋转 $90^{\circ}$ 的整数倍的变换集合构成一个群。同样，证明在二维平面上连续平移对象的变换也构成一个群（参见\cref{subsection:ch09-对称性不变性}）。
\end{exercise}

\begin{exercise}\label{exercise:ch09-02}
（★★）考虑如下形式的线性模型：
\begin{align}
y(\boldsymbol{x},\boldsymbol{w})=w_{0}+\sum_{i=1}^{D}w_{i}x_{i},
\label{align:ch09-线性模型}
\end{align}
并考虑如下形式的平方和误差函数：
\begin{align}
E_{D}(\boldsymbol{w})=\frac{1}{2}\sum_{n=1}^{N}\left\{y(\boldsymbol{x}_{n},\boldsymbol{w})-t_{n}\right\}^{2}.
\label{align:ch09-平方和误差}
\end{align}
将均值为 $0$、方差为 $\sigma^{2}$ 的 Gauss 噪声 $\varepsilon_{i}$ 独立地添加到每个输入变量 $x_{i}$ 中。利用 $\mathbb{E}[\varepsilon_{i}]=0$ 和 $\mathbb{E}[\varepsilon_{i}\varepsilon_{j}]=\delta_{ij}\sigma^{2}$，证明最小化在噪声分布上平均的 $E_{D}$ 等价于最小化无噪声输入变量的平方和误差以及权重衰减正则化项的和，其中偏置参数 $w_{0}$ 已从正则化项中省略。
\end{exercise}

\begin{exercise}\label{exercise:ch09-03}
（★★）考虑一个仅由二次正则化项组成的误差函数
\begin{align}
\Omega(\boldsymbol{w})=\frac{1}{2}\boldsymbol{w}^{\mathsf T}\boldsymbol{w}
\label{align:ch09-二次正则}
\end{align}
以及梯度下降更新公式
\begin{align}
\boldsymbol{w}^{(\tau+1)}=\boldsymbol{w}^{(\tau)}-\eta\nabla\Omega(\boldsymbol{w}).
\label{align:ch09-梯度更新}
\end{align}
通过考虑无穷小更新，写出 $\boldsymbol{w}$ 迭代的相应微分方程。从一个初值 $\boldsymbol{w}_{0}$ 开始，写出这个微分方程的解，并证明 $\boldsymbol{w}$ 的元素指数衰减到零。
\end{exercise}

\begin{exercise}\label{exercise:ch09-04}
（★）验证由\eqref{align:ch09-隐层激活}和\eqref{align:ch09-输出激活}定义的网络函数在应用于输入的转换［\eqref{align:ch09-输入线性变换}］下是不变的，前提是同时使用\eqref{align:ch09-第一层权重变换}和\eqref{align:ch09-第一层偏置变换}转换权重和偏置。同样，通过将\eqref{align:ch09-第二层权重变换}和\eqref{align:ch09-第二层偏置变换}应用到第二层权重和偏置，证明网络输出可以根据\eqref{align:ch09-输出变换}进行变换。
\end{exercise}

\begin{exercise}\label{exercise:ch09-05}
（★★）通过使用 Lagrange 乘子，证明最小化\eqref{align:ch09-lasso目标}给出的正则化误差函数等价于在约束条件［\eqref{align:ch09-lasso约束}］下最小化未正则化的误差函数 $E(\boldsymbol{w})$，并讨论参数 $\eta$ 和 $\lambda$ 之间的关系（参见\hyperref[chapter:appC-Lagrange乘子]{附录\ref*{chapter:appC-Lagrange乘子}}）。
\end{exercise}

\begin{exercise}\label{exercise:ch09-06}
（★★★）考虑如下形式的二次误差函数：
\begin{align}
E=E_{0}+\frac{1}{2}(\boldsymbol{w}-\boldsymbol{w}^{*})^{\mathsf T}\boldsymbol{H}(\boldsymbol{w}-\boldsymbol{w}^{*}),
\label{align:ch09-二次误差}
\end{align}
其中 $\boldsymbol{w}^{*}$ 代表最小值，Hessian 矩阵 $\boldsymbol{H}$ 为正定常矩阵（不随 $\boldsymbol{w}$ 变化）。假设选择初始权重向量 $\boldsymbol{w}^{(0)}$ 位于原点，并使用简单梯度下降法进行更新：
\begin{align}
\boldsymbol{w}^{(\tau)}=\boldsymbol{w}^{(\tau-1)}-\rho\nabla E,
\label{align:ch09-简单梯度下降}
\end{align}
其中 $\tau$ 表示步数，$\rho$ 表示学习率（假设学习率较小）。证明在 $\tau$ 步之后，与 $\boldsymbol{H}$ 的特征向量平行的权重向量的分量可以写成
\begin{align}
w_{j}^{(\tau)}=\left\{1-(1-\rho\eta_{j})^{\tau}\right\}w_{j}^{*},
\label{align:ch09-早停分量}
\end{align}
其中 $w_{j}=\boldsymbol{u}_{j}^{\mathsf T}\boldsymbol{w}$、$w_{j}^{*}=\boldsymbol{u}_{j}^{\mathsf T}\boldsymbol{w}^{*}$，而 $\boldsymbol{u}_{j}$ 和 $\eta_{j}$ 分别是 $\boldsymbol{H}$ 的特征向量和特征值，由下式给出：
\begin{align}
\boldsymbol{H}\boldsymbol{u}_{j}=\eta_{j}\boldsymbol{u}_{j}.
\label{align:ch09-特征方程}
\end{align}
证明当 $\tau\to\infty$ 时，有 $\boldsymbol{w}^{(\tau)}\to\boldsymbol{w}^{*}$，前提是 $|1-\rho\eta_{j}|<1$。假设训练在有限个 $\tau$ 步之后停止，证明与 $\boldsymbol{H}$ 的特征向量平行的权重向量的分量满足
\begin{align}
w_{j}^{(\tau)}\approx w_{j}^{*},\qquad\text{当 }\eta_{j}\gg(\rho\tau)^{-1}\text{ 时},
\label{align:ch09-早停大特征值}
\end{align}
\begin{align}
|w_{j}^{(\tau)}|\ll|w_{j}^{*}|,\qquad\text{当 }\eta_{j}\ll(\rho\tau)^{-1}\text{ 时}.
\label{align:ch09-早停小特征值}
\end{align}
这个结果表明 $(\rho\tau)^{-1}$ 在权重衰减中起着与正则化参数 $\lambda$ 类似的作用。
\end{exercise}

\begin{exercise}\label{exercise:ch09-07}
（★★）考虑一个其中多个权重被约束为具有相同值的神经网络。讨论在计算误差函数关于网络中可调参数的导数时，如何修改标准的反向传播算法以确保满足此类约束。
\end{exercise}

\begin{exercise}\label{exercise:ch09-08}
（★）考虑下式定义的混合分布：
\begin{align}
p(w)=\sum_{j=1}^{M}\pi_{j}\,\mathcal{N}\!\left(w\mid\mu_{j},\sigma_{j}^{2}\right),
\label{align:ch09-混合分布}
\end{align}
其中 $\pi_{j}$ 可以视为对应 Gauss 分量的先验概率 $p(j)$。使用 Bayes 定理证明相应的后验概率 $p(j\mid w)$ 由\eqref{align:ch09-后验概率}给出。
\end{exercise}

\begin{exercise}\label{exercise:ch09-09}
（★★）使用\eqref{align:ch09-Gauss混合先验}～\eqref{align:ch09-后验概率}验证\eqref{align:ch09-权重梯度}的结果。
\end{exercise}

\begin{exercise}\label{exercise:ch09-10}
（★★）使用\eqref{align:ch09-Gauss混合先验}～\eqref{align:ch09-后验概率}验证\eqref{align:ch09-均值梯度}的结果。
\end{exercise}

\begin{exercise}\label{exercise:ch09-11}
（★★）使用\eqref{align:ch09-Gauss混合先验}～\eqref{align:ch09-后验概率}验证\eqref{align:ch09-方差参数梯度}的结果。
\end{exercise}

\begin{exercise}\label{exercise:ch09-12}
（★★）证明\eqref{align:ch09-softmax混合系数}定义的混合系数 $\pi_{k}$ 相对于辅助参数 $\eta_{j}$ 的导数由下式给出：
\begin{align}
\frac{\partial\pi_{k}}{\partial\eta_{j}}=\delta_{kj}\pi_{j}-\pi_{j}\pi_{k}.
\label{align:ch09-混合系数导数}
\end{align}
因此，通过对所有 $i$ 使用约束 $\sum_{k}\gamma_{k}(w_{i})=1$，可以得到\eqref{align:ch09-混合系数梯度}的结果。
\end{exercise}

\begin{exercise}\label{exercise:ch09-13}
（★）验证\eqref{align:ch09-残差1}～\eqref{align:ch09-残差3}可以组合出一个形为\eqref{align:ch09-残差展开}的涉及整个网络的单一整体方程。
\end{exercise}

\begin{exercise}\label{exercise:ch09-14}
（★★）对于一个简单的集成模型，平均期望平方和误差 $E_{\mathrm{AV}}$ 可以用\eqref{align:ch09-平均误差}来定义，集成模型本身的期望误差由\eqref{align:ch09-集成误差}给出。假设个体误差满足\eqref{align:ch09-零均值}和\eqref{align:ch09-不相关}，请推导出\eqref{align:ch09-集成误差结果}的结果。
\end{exercise}

\begin{exercise}\label{exercise:ch09-15}
（★★）利用 Jensen 不等式针对凸函数 $f(x)=x^{2}$ 的特殊情况，证明\eqref{align:ch09-平均误差}给出的简单集成模型成员的平均期望平方和误差 $E_{\mathrm{AV}}$ 和\eqref{align:ch09-集成误差}给出的集成模型本身的期望误差 $E_{\mathrm{COM}}$ 满足
\begin{align}
E_{\mathrm{COM}}\leqslant E_{\mathrm{AV}}.
\label{align:ch09-集成上界}
\end{align}
\end{exercise}

\begin{exercise}\label{exercise:ch09-16}
（★★）利用 Jensen 不等式，证明\eqref{align:ch09-集成上界}适用于任何误差函数 $E(y)$，而非仅仅适用于平方和误差函数，只要它是 $y$ 的凸函数即可。
\end{exercise}

\begin{exercise}\label{exercise:ch09-17}
（★★）考虑一个集成模型，其允许各组成模型的权重不相等，于是有
\begin{align}
y_{\mathrm{COM}}(\boldsymbol{x})=\sum_{m=1}^{M}\alpha_{m}y_{m}(\boldsymbol{x}).
\label{align:ch09-加权集成}
\end{align}
为了确保预测 $y_{\mathrm{COM}}(\boldsymbol{x})$ 保持在合理的范围内，假设我们要求其在 $\boldsymbol{x}$ 的每个取值上都由最小值和最大值限定，这些边界可以由集成模型的任何成员给出，从而
\begin{align}
y_{\min}(\boldsymbol{x})\leqslant y_{\mathrm{COM}}(\boldsymbol{x})\leqslant y_{\max}(\boldsymbol{x}),
\label{align:ch09-加权集成界}
\end{align}
其中 $y_{\min}(\boldsymbol{x})=\min_{m}y_{m}(\boldsymbol{x})$、$y_{\max}(\boldsymbol{x})=\max_{m}y_{m}(\boldsymbol{x})$。证明这一约束的充要条件是系数 $\alpha_{m}$ 满足
\begin{align}
\alpha_{m}\geqslant0\quad\text{且}\quad\sum_{m=1}^{M}\alpha_{m}=1.
\label{align:ch09-权重条件}
\end{align}
\end{exercise}

\begin{exercise}\label{exercise:ch09-18}
（★★★）探索 dropout 正则化对使用最小二乘法训练的简单线性回归模型的影响。考虑如下形式的模型：
\begin{align}
y_{k}=\sum_{i=1}^{D}w_{ki}x_{i},
\label{align:ch09-dropout线性模型}
\end{align}
并考虑下式给出的平方和误差函数：
\begin{align}
E(\boldsymbol{w})=\sum_{n=1}^{N}\sum_{k=1}^{K}\left\{t_{nk}-\sum_{i=1}^{D}w_{ki}R_{ni}x_{ni}\right\}^{2},
\label{align:ch09-dropout误差}
\end{align}
其中 dropout 矩阵的元素 $R_{ni}\in\{0,1\}$ 是从参数 $\rho$ 的 Bernoulli 分布中随机采样得到的。对随机 dropout 参数的分布取期望。证明
\begin{align}
\mathbb{E}[R_{ni}]=\rho,
\label{align:ch09-dropout均值}
\end{align}
\begin{align}
\mathbb{E}[R_{ni}R_{nj}]=\delta_{ij}\rho+(1-\delta_{ij})\rho^{2}.
\label{align:ch09-dropout协方差}
\end{align}
然后证明 dropout 模型的期望误差函数由下式给出：
\begin{align}
\mathbb{E}[E(\boldsymbol{w})]
&=\sum_{n=1}^{N}\sum_{k=1}^{K}\left\{y_{nk}-\rho\sum_{i=1}^{D}w_{ki}x_{ni}\right\}^{2}\nonumber\\
&\qquad+\rho(1-\rho)\sum_{n=1}^{N}\sum_{k=1}^{K}\sum_{i=1}^{D}w_{ki}^{2}x_{ni}^{2}.
\label{align:ch09-dropout期望误差}
\end{align}
\enlargethispage{2\baselineskip}
可以看到，dropout 模型的期望误差函数对应于具有二次正则化项的平方和误差函数，后者的正则化系数已根据输入值分别为每个输入变量做了缩放。最后，写下用于最小化这个正则化误差函数的权重矩阵的闭式解。
\end{exercise}

\end{document}
